% open-logic.tex
% compiles to produce complete text using open-logic-dev style

\documentclass[../include/open-logic-part]{subfiles}

\begin{document}

\clearpage

\begin{editorial}
本文件加载了Open Logic Project所包含的所有内容。
像这样的编者按如果显示出来，则表明该文件在编译时并未考虑这些材料的呈现方式。
如果你能看到这段文字，那么使用此 PDF 进行教学或自学很可能是\emph{不可取的}。

Open Logic Project提供了多种机制，可以生成更适合教学或自学的文本。
例如，默认情况下，文本会将所有逻辑算子作为初始符号，并在证明中给出所有算子的所有情形。
但将其中一些情形留作练习会更好。Open Logic Project本身仍在不断建设中。
为了促进协作与改进，即使某些材料仅为草稿、缺少习题等，也将其收录其中。
供课程使用的 PDF 会将这些章节排除在外。

若要寻找更适合教学或自学的 PDF，请查看\href{http://builds.openlogicproject.org/}{OLP 网站上提供的示例课程}。
若要制作你自己的版本，可以从\href{https://github.com/OpenLogicProject/OpenLogic/tree/master/courses/sample}{示例驱动文件}入手，
或者参考衍生教材的源码，以获取更精巧、更高级的示例。
\end{editorial}

\olimport[sets-functions-relations]{sets-functions-relations-complete}

\olimport[propositional-logic]{propositional-logic}

\olimport[first-order-logic]{first-order-logic}

\olimport[model-theory]{model-theory}

\olimport[computability]{computability}

\olimport[turing-machines]{turing-machines}

\olimport[incompleteness]{incompleteness}

\olimport[second-order-logic]{second-order-logic}

\olimport[lambda-calculus]{lambda-calculus}

\olimport[many-valued-logic]{many-valued-logic}

\olimport[normal-modal-logic]{normal-modal-logic}

\olimport[applied-modal-logic]{applied-modal-logic}

\olimport[intuitionistic-logic]{intuitionistic-logic}

\olimport[counterfactuals]{counterfactuals}

\olimport[set-theory]{set-theory}

\olimport[methods]{methods}

\olimport[history]{history}

\olimport[reference]{reference}

\end{document}